
\emph{A Euclidean Gibbs probe measure with automatic first moment.} This paragraph records a conditional construction of a measure with finite one-point potential moment, supplying the single consumable estimate left open by the paragraph above, and its typed one-point consequences. Let \(\mathrm{dvol}\) be the volume element of the Presentation-B background, assumed closed with finite volume \(\mathrm{Vol}\), fix the background gauge field of Equation~\eqref{eq:Smicro}, and take as hypothesis its Euclidean continuation \(S_E[\Phi]=\int(\lvert D\Phi\rvert_F^2+\mathcal V(\Phi))\,\mathrm{dvol}\) with \(\mathcal V\) the potential of Equation~\eqref{eq:V}. Restrict \(\Phi\) to a finite-dimensional real-linear space of matrix fields containing the constant fields, with coordinates \(\phi\in\mathbb R^K\) and \(\phi\mapsto\Phi(x)\) linear and non-zero at each \(x\) (the truncation), and define the Gibbs history measure \(\mu=\mathcal Z^{-1}e^{-S_E}\prod_k\mathrm d\phi_k\) with normalizing integral \(\mathcal Z\), keeping the history measure \(\mu\) notationally distinct from the volume element \(\mathrm{dvol}\). Then: (i) \(0<\mathcal Z<\infty\): \(S_E(\phi)\to\infty\) as \(\lVert\phi\rVert_2\to\infty\) (since \(\int\lvert\Phi\rvert_F^4\,\mathrm{dvol}=0\) forces \(\Phi\equiv0\) and the pointwise coercivity \(\mathcal V\ge(\lambda/32)\lVert\Phi\rVert_F^4-2\lambda v^4\), verified below, controls the quartic uniformly on the unit sphere of the truncation), so \(\mathcal Z\) is dominated by a super-Gaussian integral; and \(\mathcal Z>0\) because \(S_E\) is bounded on the tube of constant fields \(v(\mathbb I+\{\lVert E\rVert_F\le\rho\})\) about the vacuum manifold Equation~\eqref{eq:vacuum-manifold}. (ii) The first moment is automatic: the pointwise interpolation \(\sup_x\lVert\Phi(x)\rVert_F\le C\lVert\phi\rVert_2\) on the truncation, Jensen's inequality, and the coercivity bounds give \(\mathcal V(x)\le(8C^4/\mathrm{Vol})\int\mathcal V\,\mathrm{dvol}+\lambda v^4(16C^4+1)\le(8C^4/\mathrm{Vol})\,S_E+\lambda v^4(16C^4+1)\), whence \(\mathbb E_\mu[\mathcal V(x)]\le(8C^4/\mathrm{Vol})\,\mathbb E_\mu[S_E]+\lambda v^4(16C^4+1)\) and, since \(\sup_{s\ge0}s\,e^{-s/2}=2/e\), \(\mathbb E_\mu[S_E]\le(2/e)\,\mathcal Z(1/2)/\mathcal Z(1)\), where \(\mathcal Z(\beta)\) is the normalizing integral of \(e^{-\beta S_E}\). (iii) The estimate of the paragraph above, applied to the one-point distribution \(\omega_x\) of \(\Phi(x)\), gives \(\omega_x\{\det e(x)=0\}\le\omega_x\{\mathcal V(x)\ge\lambda v^4/4\}\le\min\{1,4\,\mathbb E_{\omega_x}[\mathcal V(x)]/(\lambda v^4)\}\) under the same configuration-level reading of \(e\) through Equation~\eqref{eq:vev-is-frame}; and on this measure the event itself has zero mass, \(\omega_x\{\det e(x)=0\}=0\) identically: the singular set is a proper real-algebraic cone and the constant identity field lies in the truncation, so its inverse image under \(\phi\mapsto\Phi(x)\) is a proper algebraic subset of \(\mathbb R^K\), which the \(e^{-S_E}\)-density cannot charge. (iv) By Fubini's theorem, \(\mu\)-almost every configuration has spacetime-Lebesgue-null defect set \(\{x:\det\Phi(x)=0\}\) - null, not empty, so this does not discharge the pointwise Assumption~\ref{ass:nondeg}. (v) For the inverse-temperature family \(\mu_\beta\propto e^{-\beta S_E}\), convexity of \(\ln\mathcal Z(\beta)\) makes \(\mathbb E_{\mu_\beta}[S_E]\) nonincreasing in \(\beta\), so the volume-averaged soft-event estimate \(\mathrm{Vol}^{-1}\int\omega_x\{\mathcal V(x)\ge\lambda v^4/4\}\,\mathrm{dvol}\le4\,\mathbb E_{\mu_\beta}[S_E]/(\lambda v^4)\le4\,\mathbb E_{\mu_{\beta_0}}[S_E]/(\lambda v^4)\) holds for every \(\beta\ge\beta_0\). Scope: the hypotheses are stated and the conclusions are conditional derivations within them; \(\mu\) is a Euclidean Gibbs probe measure on a truncated Presentation-B field theory and not a state of the matrix theory, its Wick rotation, the closedness of the background and the choice of truncation are assumed rather than derived, the estimate lives at configuration level so the vacuum-expectation reading \(e=\langle\Phi\rangle/v\) of Equation~\eqref{eq:vev-is-frame} would additionally require control of the mean field \(\mathbb E_\mu[\Phi(x)]\) that is not provided here, removal of the truncation is open, and the selection of Equation~\eqref{eq:V} by Presentation~A (Section~\ref{sec:ch3-relate}) remains the upstream input that gates what any Presentation-B measure can be cited as; Assumptions~\ref{ass:frame-lock} and~\ref{ass:nondeg} remain assumed, so Problem~D remains OPEN and PROGRAMME-CONDITIONAL.
